% ===== Appendix: corrections to the book, what we added, sources (after Questions; not presented) =====
\appendix
\partaccent{lavender}\gdef\idiottopicemoji{books}\gdef\idiottopicname{Հավելված}
\newcolumntype{L}[1]{>{\raggedright\arraybackslash}p{#1}}

\begin{frame}{Հավելված Ա. Ուղղումներ գրքի նկատմամբ}
  \footnotesize
  \setlength{\tabcolsep}{4pt}
  \renewcommand{\arraystretch}{1.35}
  \begin{tabular}{@{}L{2.4cm}L{3.5cm}L{4.45cm}L{3.2cm}@{}}
    \toprule
    \textbf{Թեմա} & \textbf{Գրքում (2016)} & \textbf{Այսօր} & \textbf{Աղբյուր}\\
    \midrule
    Սուզորդների փորձը & Նույն միջավայրը լավացնում է վերհիշումը & 16 սուզորդով կրկնությունը էֆեկտը չգտավ & Murre, R. Soc. Open Sci., 2021\\
    Պացիենտ Հ.Մ. & Մահացել է 78 տարեկանում & 82 տարեկանում (1926-2008) & Squire, Neuron, 2009\\
    \bottomrule
  \end{tabular}
\end{frame}

\begin{frame}{Հավելված Բ. Ինչ ենք ավելացրել գրքից դուրս}
  \begin{columns}[T]
    \begin{column}{0.48\textwidth}
      \begin{card}[colback=mint!30!cardwhite, fontupper=\scriptsize]
        {\blackfont Թվեր ու հետազոտություններ}\par\smallskip
        \begin{itemize}\setlength{\itemsep}{0.5mm}
          \item «Քուն» բառի թեստը (DRM)՝ Roediger \& McDermott, 1995
          \item Լոֆթուսի թվերը՝ Loftus \& Palmer, 1974
          \item Տաքսու վարորդների հիպոկամպը՝ Maguire et al., 2000
        \end{itemize}
      \end{card}
    \end{column}
    \begin{column}{0.48\textwidth}
      \begin{card}[colback=peach!35!cardwhite, fontupper=\scriptsize]
        {\blackfont Մերն է, ոչ թե գրքից}\par\smallskip
        \begin{itemize}\setlength{\itemsep}{0.5mm}
          \item Անուն հիշելու 2-րդ ու 3-րդ խորհուրդները
          \item Թվերի խմբավորման (chunking) օրինակը
          \item Հ.Մ.-ի հայելու փորձի մանրամասները (Բրենդա Միլներ). գրքում՝ «նկարչական հմտություններ»
          \item Մեմը. բովանդակությունը գրքից է
        \end{itemize}
      \end{card}
    \end{column}
  \end{columns}
\end{frame}

\begin{frame}{Հավելված Գ. Աղբյուրներ. գլուխ 2}
  \scriptsize
  Գրքի ծանոթագրություններից, որոնց վրա հիմնված են սլայդները\par\smallskip
  \begin{columns}[T]
    \begin{column}{0.48\textwidth}
      \begin{itemize}\setlength{\itemsep}{0.4mm}
        \item Cowan N. The magical mystery four. Curr. Dir. Psychol. Sci., 2010
        \item Miller G. A. The magical number seven, plus or minus two. Psychol. Rev., 1956 (գիրքը հղում է Nicolis \& Tsuda, 1985-ին)
        \item Godden D. R., Baddeley A. D. Context-dependent memory in two natural environments. Br. J. Psychol., 1975
        \item Henson R. N. Short-term memory for serial order. Cogn. Psychol., 1998
        \item Eichenbaum H. The Cognitive Neuroscience of Memory. 2011
        \item Greenwald A. G. The totalitarian ego. Am. Psychol., 1980
        \item Neisser U. John Dean's memory: a case study. Cognition, 1981
      \end{itemize}
    \end{column}
    \begin{column}{0.48\textwidth}
      \begin{itemize}\setlength{\itemsep}{0.4mm}
        \item Mather M., Johnson M. K. Choice-supportive source monitoring. Psychol. Aging, 2000
        \item Meissner C. A., Brigham J. C. Thirty years of investigating the own-race bias. Psychol. Public Policy Law, 2001
        \item Hoffrage U. et al. Hindsight bias: a by-product of knowledge updating? JEP:LMC, 2000
        \item Walker W. R., Skowronski J. J. The fading affect bias. Appl. Cogn. Psychol., 2009
        \item Wilson A. E., Ross M. From chump to champ. JPSP, 2001
        \item Loftus E. F. The price of bad memories. 1998 (Նադին Քուլի դեպքը)
        \item Squire L. R. The legacy of patient H.M. for neuroscience. Neuron, 2009
      \end{itemize}
    \end{column}
  \end{columns}
\end{frame}
